\section{Proof of Theorem \ref{thm:perturb}}

\subsection{The disk inverse step}\label{subsec:disk-inverse-step}

The proof rests on a local inverse theorem at the Daubechies disk. The result
below is self-contained, but its point of view is inspired by perturbative
inverse-source and quadrature-domain methods: Isakov's local free-boundary
scheme for inverse source problems, Cherednichenko's local inverse logarithmic
potential theorem, and the quadrature-domain framework developed by
Gustafsson, Shahgholian, Karp, and Aharonov--Schiffer--Zalcman. In the present
disk-based setting the required inverse step can be proved directly by
diagonalizing the full moment map in Fourier modes.

Fix \(0<\lambda<1\), and let \(r_\lambda>0\) be determined by
\[
1-e^{-\pi r_\lambda^2}=\lambda .
\]
For a real-valued \(2\pi\)-periodic function \(\varphi\) with
\(\|\varphi\|_{C^0}<r_\lambda/2\), set
\[
U_\varphi
:=
\{re^{i\theta}:0\le r<r_\lambda+\varphi(\theta)\}.
\]
For \(\tau>0\), let \(\mathcal A_\tau^\R\) be the Banach space of real-valued
periodic functions
\[
\varphi(\theta)=\sum_{k\in\Z}\varphi_k e^{ik\theta},
\qquad
\varphi_{-k}=\overline{\varphi_k},
\]
with norm
\[
\|\varphi\|_{\mathcal A_\tau}
:=
\sum_{k\in\Z}|\varphi_k|e^{\tau |k|}<\infty .
\]
This is a Banach algebra, and \(\mathcal A_\tau^\R\hookrightarrow C^m\) for
every fixed \(m\). In particular, small elements of \(\mathcal A_\tau^\R\)
produce real-analytic \(C^2\) Jordan domains.

\begin{proposition}[Disk symmetrized inverse theorem]\label{prop:disk-inverse-tailored}
Let \(N\ge1\), \(0<\lambda<1\), and \(r_\lambda\) be as above. There are
\(\tau>0\), \(\varepsilon>0\), and a \(C^1\) map
\[
b=(b_1,\ldots,b_N)\longmapsto \varphi_b\in\mathcal A_\tau^\R,
\qquad
\varphi_0=0,
\]
defined for \(\max_j|b_j|<\varepsilon\), such that, for
\[
P_b(z):=1+\sum_{j=1}^N b_jz^j,
\qquad
P_b^\#(z):=\overline{P_b(\overline z)},
\qquad
Q_b(w):=P_b^\#(\partial_w/\pi)P_b(w),
\]
and
\[
U_b:=U_{\varphi_b},
\]
one has
\begin{equation}\label{eq:disk-inverse-identity}
\int_{U_b}|P_b(z)|^2e^{-\pi |z|^2}e^{\pi\overline z w}\,dA(z)
=
\lambda Q_b(w),
\qquad w\in\C .
\end{equation}
Moreover,
\[
\|\varphi_b\|_{C^2(\mathbb S^1)}
\le C_{\lambda,N}\max_{1\le j\le N}|b_j|,
\]
and the branch is locally unique among sufficiently small analytic radial
graphs in \(\mathcal A_\tau^\R\).
\end{proposition}

\begin{proof}
We solve the coefficient equations obtained by expanding the entire identity
\eqref{eq:disk-inverse-identity}. Write
\[
Q_b(w)=\sum_{j=0}^N q_j(b)w^j,
\]
and \(q_j(b)=0\) for \(j>N\). The constant coefficient is real:
\[
q_0(b)=1+\sum_{k=1}^N \frac{k!}{\pi^k}|b_k|^2.
\]
Since
\[
e^{\pi\overline z w}
=
\sum_{j=0}^\infty \frac{\pi^j}{j!}\overline z^{\,j}w^j,
\]
the desired identity is equivalent to the infinite system
\begin{equation}\label{eq:disk-moment-system}
\frac{\pi^j}{j!}\int_{U_\varphi}
|P_b(z)|^2\overline z^{\,j}e^{-\pi |z|^2}\,dA(z)
=
\lambda q_j(b),
\qquad j\ge0 .
\end{equation}

We introduce the weighted sequence space adapted to the linearization at the
circle. Put
\[
a_j:=2\pi e^{-\pi r_\lambda^2}r_\lambda^{j+1}\frac{\pi^j}{j!},
\qquad j\ge0,
\]
and let \(\mathcal Y_\tau\) be the space of sequences
\(y=(y_j)_{j\ge0}\), with \(y_0\in\R\), such that
\[
\|y\|_{\mathcal Y_\tau}
:=
\frac{|y_0|}{a_0}
+2\sum_{j=1}^\infty e^{\tau j}\frac{|y_j|}{a_j}
<\infty .
\]
Define
\[
\mathcal F_j(b,\varphi)
:=
\frac{\pi^j}{j!}\int_{U_\varphi}
|P_b(z)|^2\overline z^{\,j}e^{-\pi |z|^2}\,dA(z)
-\lambda q_j(b),
\qquad j\ge0 .
\]
Thus \(\mathcal F(b,\varphi)=0\) is precisely
\eqref{eq:disk-moment-system}. At \(b=0\), \(\varphi=0\), the domain is the
disk \(D_{r_\lambda}\), and radial symmetry gives
\[
\frac{\pi^j}{j!}\int_{D_{r_\lambda}}
\overline z^{\,j}e^{-\pi |z|^2}\,dA(z)
=
\begin{cases}
\lambda, & j=0,\\
0, & j\ge1.
\end{cases}
\]
Hence \(\mathcal F(0,0)=0\).

We next compute the derivative with respect to the boundary graph. Let
\(\psi(\theta)=\sum_{k\in\Z}\psi_k e^{ik\theta}\in\mathcal A_\tau^\R\).
Differentiating the polar-coordinate formula
\[
\int_{U_\varphi}H(z)\,dA(z)
=
\int_0^{2\pi}\int_0^{r_\lambda+\varphi(\theta)}
H(re^{i\theta})r\,dr\,d\theta
\]
at \(\varphi=0\), with
\(H(z)=\overline z^{\,j}e^{-\pi|z|^2}\), gives
\[
D_\varphi\mathcal F_j(0,0)[\psi]
=
\frac{\pi^j}{j!}
\int_0^{2\pi}
\psi(\theta)r_\lambda^{j+1}e^{-ij\theta}e^{-\pi r_\lambda^2}
\,d\theta
=
a_j\psi_j .
\]
Consequently
\[
D_\varphi\mathcal F(0,0)\colon
\mathcal A_\tau^\R\to\mathcal Y_\tau,
\qquad
\psi\mapsto(a_j\psi_j)_{j\ge0},
\]
is an isomorphism. Its inverse is explicit:
if \(y\in\mathcal Y_\tau\), then
\[
\psi_0=\frac{y_0}{a_0},
\qquad
\psi_j=\frac{y_j}{a_j}\quad(j\ge1),
\qquad
\psi_{-j}=\overline{\psi_j}.
\]
The definition of the two norms gives
\(\|\psi\|_{\mathcal A_\tau}=\|y\|_{\mathcal Y_\tau}\).

It remains to justify the use of the Banach-space implicit function theorem.
For \(\|\varphi\|_{\mathcal A_\tau}\) small, the function
\(\theta\mapsto r_\lambda+\varphi(\theta)\) extends holomorphically to the
strip \(|\Im\theta|<\tau\) and stays in a fixed annulus
\(r_\lambda/2<|r|<3r_\lambda/2\). In polar coordinates,
\[
\mathcal F_j(b,\varphi)+\lambda q_j(b)
=
\frac{\pi^j}{j!}\int_0^{2\pi}e^{-ij\theta}
\int_0^{r_\lambda+\varphi(\theta)}
P_b(re^{i\theta})P_b^\#(re^{-i\theta})r^{j+1}e^{-\pi r^2}\,dr\,d\theta .
\]
Expanding the inner integral at \(r_\lambda\) gives an absolutely convergent
power series in \(\varphi\), whose coefficients are analytic functions of
\(\theta\) and polynomial functions of \(b\). Since \(\mathcal A_\tau\) is a
Banach algebra and \(P_b(re^{i\theta})P_b^\#(re^{-i\theta})\) is a finite
trigonometric polynomial of degree at most \(N\), the resulting series and
its first derivatives in \((b,\varphi)\) converge in \(\mathcal Y_\tau\) after
possibly decreasing \(\tau\) and restricting to
\(\|b\|+\|\varphi\|_{\mathcal A_\tau}<\eta\). More explicitly, Cauchy's
estimate on a slightly wider strip gives, for \(0<\tau'<\tau\),
\[
\sum_{j\ge0}e^{\tau' j}\frac1{a_j}
\left|
\frac{\pi^j}{j!}\widehat{R_j}(j)
\right|
\le
C_{\lambda,N,\tau,\tau'}\|R\|_{\mathcal A_\tau},
\]
where \(\widehat{R_j}(j)\) denotes the \(j\)-th Fourier coefficient of the
analytic coefficient appearing in the \(j\)-th moment equation. Applying this
estimate term by term to the Taylor expansion in \(\varphi\) yields
\[
\|\mathcal F(b,\varphi)-\mathcal F(b,\psi)
-D_\varphi\mathcal F(b,\psi)[\varphi-\psi]\|_{\mathcal Y_{\tau'}}
\le
C(\|b\|+\|\varphi\|_{\mathcal A_\tau}
+\|\psi\|_{\mathcal A_\tau})
\|\varphi-\psi\|_{\mathcal A_\tau}^2,
\]
and analogous estimates for differentiation in \(b\). Shrinking the strip
once and then relabelling \(\tau'\) as \(\tau\), this proves that
\[
\mathcal F:
\{(b,\varphi):\|b\|+\|\varphi\|_{\mathcal A_\tau}<\eta\}
\longrightarrow \mathcal Y_\tau
\]
is \(C^1\).

The Banach implicit function theorem now applies at \((0,0)\), because
\(D_\varphi\mathcal F(0,0)\) is an isomorphism. Thus, for \(\|b\|\) small,
there is a unique \(\varphi_b\in\mathcal A_\tau^\R\), small in
\(\mathcal A_\tau\), such that \(\mathcal F(b,\varphi_b)=0\), and
\(b\mapsto\varphi_b\) is \(C^1\). The embedding
\(\mathcal A_\tau\hookrightarrow C^2\) gives the displayed \(C^2\) bound.
Finally, \(\mathcal F(b,\varphi_b)=0\) gives all coefficient identities
\eqref{eq:disk-moment-system}; summing the power series in \(w\) gives
\eqref{eq:disk-inverse-identity} for every \(w\in\C\), since both sides are
entire functions. This proves the proposition.
\end{proof}

\begin{remark}
The proof above is the only place where analyticity of the boundary
parametrization is used. The constructed domains are therefore \(C^\infty\),
hence \(C^2\), which is the regularity needed later. The argument does not
assert a general local quadrature-domain theorem near an arbitrary reference
domain; it proves exactly the disk-based inverse statement used below.
\end{remark}

\begin{proof}[Proof of Theorem \ref{thm:perturb}]
We again use the Bargmann transform. Let
\[
P(z)=1+\sum_{k=1}^N b_k z^k,
\]
where the coefficients $b_k$ are chosen so that the Bargmann transform of
\[
f_0=h_0+\sum_{k=1}^N a_k h_k
\]
is exactly $P$. We also write
\[
P^\#(z):=\overline{P(\overline z)}
=1+\sum_{k=1}^N \overline{b_k}z^k
\]
for the coefficientwise conjugate polynomial. Then the condition that $f_0$ be an eigenfunction of $\mathcal{L}_U$ with eigenvalue $\lambda$ is equivalent to
\begin{equation}\label{eq:eigenfunction-general-new}
\int_U P(z)e^{-\pi|z|^2}e^{\pi \overline z w}\,dA(z)=\lambda P(w),
\qquad w\in \C.
\end{equation}
Set
\[
Q(w):=P^\#(\partial_w/\pi)P(w).
\]
By Proposition~\ref{prop:disk-inverse-tailored}, after decreasing
\(\varepsilon_0(\lambda,N)\) if necessary, there exists a real-analytic radial
graph \(U=U_b\), \(C^2\)-close to the Daubechies disk of radius
\(r_\lambda\), for which
\begin{equation}\label{eq:almost-reduced-new}
\int_U |P(z)|^2e^{-\pi|z|^2}e^{\pi\overline z w}\,dA(z)
=
\lambda Q(w),
\qquad w\in\C .
\end{equation}
We now recover the unsymmetrized identity
\eqref{eq:eigenfunction-general-new}. Define
\[
G(w):=\int_U P(z)e^{-\pi|z|^2}e^{\pi\overline z w}\,dA(z).
\]
The function \(G\) is entire, and differentiating under the integral is
justified by boundedness of \(U\). Since
\[
\Big(\frac{\partial_w}{\pi}\Big)^j e^{\pi\overline z w}
=\overline z^{\,j}e^{\pi\overline z w},
\]
we have
\[
P^\#(\partial_w/\pi)G(w)
=
\int_U |P(z)|^2e^{-\pi|z|^2}e^{\pi\overline z w}\,dA(z)
=
\lambda Q(w)
=
P^\#(\partial_w/\pi)(\lambda P)(w).
\]
Thus \(H:=G-\lambda P\) is an entire solution of
\[
P^\#(\partial_w/\pi)H=0.
\]

Let \(\widetilde P^\#(\zeta):=P^\#(\zeta/\pi)\), and let
\(\zeta_1,\ldots,\zeta_\ell\) be its zeros, with multiplicities
\(m_1,\ldots,m_\ell\). The elementary structure theorem for
constant-coefficient ordinary differential equations gives
\[
H(w)=\sum_{j=1}^\ell R_j(w)e^{\zeta_j w},
\]
where \(\deg R_j<m_j\). On the other hand, if
\(M:=\sup_{z\in U}|z|\), then
\[
|G(w)|\le C(1+|w|)^N e^{\pi M|w|},
\qquad w\in\C,
\]
and therefore \(H\) has exponential type at most \(\pi M\).
Since \(U\) stays in a fixed neighborhood of \(D_{r_\lambda}\), we have
\(M\le C_\lambda\). The zeros of \(P^\#\), and hence of
\(\widetilde P^\#\), escape to infinity as \(b\to0\). Shrinking
\(\varepsilon_0(\lambda,N)\) once more, we may assume
\[
|\zeta_j|>10M,\qquad 1\le j\le \ell .
\]
If some \(R_j\) were nonzero, choose one of maximal \(|\zeta_j|\) for which
\(R_j\not\equiv0\) and relabel it \(\zeta_1\). Since the \(\zeta_j\) are
\emph{distinct} (they are the distinct roots of \(\widetilde P^\#\)),
the Cauchy--Schwarz inequality
\(\Re(\zeta_j\overline{\zeta_1})\le|\zeta_j||\zeta_1|\) is strict whenever
\(\zeta_j\ne\zeta_1\) and \(|\zeta_j|=|\zeta_1|\); when \(|\zeta_j|<|\zeta_1|\)
strictness is automatic. Hence for every \(\zeta_j\ne\zeta_1\),
\[
\Re(\zeta_j\overline{\zeta_1})<|\zeta_1|^2 .
\]
Along \(w=t\overline{\zeta_1}\), \(t\to+\infty\), the term
\(R_1(w)e^{\zeta_1w}\) has modulus \(|R_1(t\overline{\zeta_1})|e^{t|\zeta_1|^2}\),
which grows like a positive power of \(t\) times \(e^{t|\zeta_1|^2}\)
because \(R_1\not\equiv0\); every other term is bounded by a polynomial
in \(t\) times \(e^{t\Re(\zeta_j\overline{\zeta_1})}\), with strictly smaller
exponential rate. Therefore
\[
|H(t\overline{\zeta_1})|\ge c e^{t|\zeta_1|^2}
\]
for all large \(t\). The exponential-type bound gives, for every
\(\delta>0\),
\[
|H(t\overline{\zeta_1})|
\le C_\delta e^{(\pi+\delta)Mt|\zeta_1|}.
\]
Letting \(t\to\infty\) and then \(\delta\downarrow0\) yields
\(|\zeta_1|\le \pi M\), contradicting \(|\zeta_1|>10M\). Hence all
\(R_j\) vanish and \(H\equiv0\), which proves
\eqref{eq:eigenfunction-general-new}.

The Bargmann transform is unitary and intertwines the localization operator
with the integral operator in \eqref{eq:eigenfunction-general-new}. Therefore
the inverse Bargmann transform \(f_0\) of \(P\) satisfies
\[
\mathcal L_U f_0=\lambda f_0 .
\]
This proves the theorem.
\end{proof}

\iffalse
The proof proceeds in four steps. In Step~0 we apply the coefficientwise conjugate differential operator $P^\#(\partial_w/\pi)$ to \eqref{eq:eigenfunction-general-new} in order to symmetrize the integrand and reduce matters to a moment identity for $|P|^2 e^{-\pi|z|^2}$. In Step~1 we construct an auxiliary radially symmetric domain $U_0$ for which the corresponding eigenvalue is the constant function $\lambda$; this is achieved by perturbing a Daubechies disk via the Isakov free-boundary scheme of Proposition~\ref{prop:isakov-tailored}. In Step~2 we use a Laplace transform in the variable $w$ to convert the polynomial right-hand side $\lambda Q(w)$ into the rational right-hand side required by Proposition~\ref{prop:cherednichenko-tailored} and produce a domain $U$ satisfying the symmetrized moment identity. Finally, in Step~3 we invert the differentiation of Step~0 by analyzing the entire solutions of the constant-coefficient ODE $P^\#(\partial_w/\pi)H=0$ and ruling out nontrivial exponential solutions through a growth estimate.

\medskip
\noindent\textbf{Step 0: a first reduction.}
For a polynomial
\[
R(z)=\alpha_0+\alpha_1 z+\cdots+\alpha_n z^n,
\]
define the differential operator
\[
R(\partial_w/\pi):=\sum_{k=0}^n \alpha_k\Big(\frac{\partial_w}{\pi}\Big)^k.
\]
We apply $P^\#(\partial_w/\pi)$ to both sides of \eqref{eq:eigenfunction-general-new}. The right-hand side is the polynomial $\lambda P(w)$, on which the differential operator acts in the obvious way. For the left-hand side, write
\[
F(w):=\int_U P(z)e^{-\pi|z|^2}e^{\pi \overline z w}\,dA(z).
\]
Since $U$ is bounded, the integrand is jointly continuous on $\overline U\times \C$, and for every compact set $K\subset\C$ and every nonnegative integer $j$ one has the uniform bound
\[
\big|\partial_w^{\,j}\big(P(z)e^{-\pi|z|^2}e^{\pi\overline z w}\big)\big|
\le
\big(\pi\,\sup_{z\in U}|z|\big)^{j}\,\sup_{z\in U}\!\big(|P(z)|e^{-\pi|z|^2}\big)\,e^{\pi(\sup_{z\in U}|z|)\sup_{w\in K}|w|},
\]
valid for $(z,w)\in U\times K$. The right-hand side is integrable on the bounded set $U$, so dominated convergence justifies differentiating under the integral sign arbitrarily often, and $F$ is entire. Combining this with the identity
\[
\Big(\frac{\partial_w}{\pi}\Big)^j e^{\pi \overline z w}=\overline z^{\,j} e^{\pi \overline z w},
\]
we obtain
\[
P^\#(\partial_w/\pi)F(w)
=
\int_U P(z)\,\overline{P(z)}\,e^{-\pi|z|^2}e^{\pi\overline z w}\,dA(z)
=
\int_U |P(z)|^2 e^{-\pi|z|^2}e^{\pi\overline z w}\,dA(z).
\]
Hence \eqref{eq:eigenfunction-general-new} is equivalent to
\begin{equation}\label{eq:almost-reduced-new}
\int_U |P(z)|^2 e^{-\pi|z|^2} e^{\pi \overline z w}\,dA(z)=\lambda Q(w),
\qquad w\in\C,
\end{equation}
where
\[
Q(w):=P^\#(\partial_w/\pi)\big(P(w)\big).
\]
Strictly speaking, applying $P^\#(\partial_w/\pi)$ only gives the implication \eqref{eq:eigenfunction-general-new}~$\Rightarrow$~\eqref{eq:almost-reduced-new}; the converse implication will be recovered in Step~3 by an ODE-theoretic argument. Our first task is therefore to construct a domain $U$ satisfying \eqref{eq:almost-reduced-new}.

\medskip
\noindent\textbf{Step 1: construction of an auxiliary domain.}
We begin by constructing a $C^2$ domain $U_0$, close to a fixed disk, for which
\begin{equation}\label{eq:first-domain-new}
\int_{U_0} |P(z)|^2 e^{-\pi|z|^2} e^{\pi \overline z w}\,dA(z)=\lambda,
\qquad w\in\C.
\end{equation}
Note that \eqref{eq:first-domain-new} is the special case of \eqref{eq:almost-reduced-new} in which the right-hand side is the constant $\lambda$; it will serve as the unperturbed reference for Step~2.

Let $\psi\in C_c^\infty(\C)$ be a nonnegative radial function such that
\[
\supp \psi \subset D_2(0)\setminus D_1(0),
\qquad
\int_\C \psi\,dA=1,
\]
and set
\[
\psi_\delta(z):=\delta^{-2}\psi(z/\delta),\qquad 0<\delta<1.
\]
Then $\psi_\delta$ has total mass one and $\supp\psi_\delta\subset D_{2\delta}(0)\setminus D_\delta(0)$. Since $\psi_\delta$ is radial, expanding $e^{\pi\overline z w}=\sum_{j\ge 0}(\pi\overline z w)^j/j!$ and integrating in polar coordinates kills every angular harmonic with $j\ge 1$, so that
\[
\int_{\C} \psi_\delta(z)e^{-\pi|z|^2}e^{\pi \overline z w}\,dA(z)=c_\delta,
\qquad w\in\C,
\]
where
\[
c_\delta:=\int_{\C}\psi_\delta(z)e^{-\pi|z|^2}\,dA(z)>0.
\]
Note that $c_\delta\to 1$ as $\delta\to 0$. Hence \eqref{eq:first-domain-new} is equivalent to
\begin{equation}\label{eq:second-domain-new}
\int_{U_0} |P(z)|^2 e^{-\pi|z|^2} e^{\pi \overline z w}\,dA(z)
=
\frac{\lambda}{c_\delta}
\int_{\C}\psi_\delta(z)e^{-\pi|z|^2}e^{\pi \overline z w}\,dA(z).
\end{equation}

We shall construct a compactly supported distributional solution of the Poisson equation
\begin{equation}\label{eq:inverse-first-new}
\Delta u_0
=
|P(z)|^2 e^{-\pi|z|^2}\mathbf{1}_{U_0}
-\frac{\lambda}{c_\delta}\psi_\delta(z)e^{-\pi|z|^2}.
\end{equation}
Once such a solution has been constructed, the moment identity \eqref{eq:first-domain-new} follows by applying the same Fourier-analytic divisibility argument used in the proof of Theorem~\ref{thm:eig-loc-U} to the corresponding compactly supported distribution
\[
\omega_0:=
|P(z)|^2e^{-\pi|z|^2}\mathbf 1_{U_0}
-\frac{\lambda}{c_\delta}\psi_\delta(z)e^{-\pi|z|^2}.
\]
Indeed, \eqref{eq:inverse-first-new} says exactly that $\Delta u_0=\omega_0$, and since $u_0$ is compactly supported, the vanishing of $\widehat{\omega_0}$ on the two characteristic lines of the Bargmann transform is equivalent to \eqref{eq:first-domain-new}.

\smallskip
\emph{Unperturbed reference.} We first treat the unperturbed case $P\equiv 1$. By the theorem of Daubechies \cite{Daubechies}, there exists a disk $D_0$ centered at the origin such that
\[
\int_{D_0} e^{-\pi|z|^2}e^{\pi \overline z w}\,dA(z)=\lambda,
\qquad w\in\C.
\]
Equivalently, $\Delta u_*=e^{-\pi|z|^2}\mathbf 1_{D_0}-\lambda\,\delta_0$ admits a compactly supported radial solution $u_*$. This solution is unique among compactly supported distributional solutions: by radial symmetry, the Poisson equation reduces to the ODE
\[
\frac{1}{r}\big(r u_*'(r)\big)'=e^{-\pi r^2}\mathbf 1_{(0,r_0)}(r),\qquad r>0,
\]
together with the asymptotic condition $u_*'(r)\sim -\lambda/(2\pi r)$ as $r\to 0^+$ enforced by the distributional source \(-\lambda\delta_0\). Integrating twice with the natural boundary condition that $u_*$ vanish for $r\ge r_0$ gives a unique radial profile. Explicitly,
\[
u_*(r)=\frac{\lambda}{2\pi}\log\frac{r_0}{r}-\int_r^{r_0}\frac{1}{s}\int_0^{s}\rho\,e^{-\pi\rho^2}\,d\rho\,ds,
\qquad 0<r\le r_0,
\]
and $u_*\equiv 0$ for $r\ge r_0$, where $r_0$ is the radius of $D_0$. From this explicit formula one verifies directly:
\begin{itemize}
\item $\Delta u_*=e^{-\pi|z|^2}>0$ on $D_0\setminus\{0\}$;
\item $u_*=0$ on $\partial D_0$;
\item the logarithmic singularity created by the mass \(-\lambda\delta_0\) has positive coefficient $\lambda/(2\pi)$ in \(\log(r_0/r)\), so $u_*(r)\to+\infty$ as $r\to 0^+$.
\end{itemize}
Equivalently, if \(A(s):=\int_0^s \rho e^{-\pi\rho^2}\,d\rho\), then the
Daubechies identity gives \(A(r_0)=\lambda/(2\pi)\), and the displayed formula
can be rewritten as
\[
u_*(r)=\int_r^{r_0}\frac{A(r_0)-A(s)}{s}\,ds>0,
\qquad 0<r<r_0.
\]
In particular \(u_*>0\) on every compact subset of \(D_0\setminus\{0\}\).

For $\delta$ small enough that $\supp\psi_\delta\subset D_0$, the same argument yields a compactly supported radial solution $\widetilde u_0$ of
\begin{equation}\label{eq:unperturbed-free-boundary-new}
\Delta \widetilde u_0
=
e^{-\pi|z|^2}\mathbf{1}_{D_0}
-\frac{\lambda}{c_\delta}\psi_\delta(z)e^{-\pi|z|^2},
\end{equation}
and $\widetilde u_0$ is uniquely determined among compactly supported radial distributional solutions by the same ODE reduction (the smooth radial source on the right-hand side replaces the Dirac mass, and the requirements of compact support and continuity through the origin pin down the integration constants).

We claim that, for $\delta$ sufficiently small,
\[
\widetilde u_0>0 \qquad \text{in } D_0.
\]
To see this, we first verify that
\begin{equation}\label{eq:psi-delta-to-mass}
\frac{\lambda}{c_\delta}\,\psi_\delta(z)\,e^{-\pi|z|^2}\xrightarrow[\delta\to 0]{}\lambda\,\delta_0
\qquad \text{in }\mathcal D'(\C).
\end{equation}
Indeed, for any $\varphi\in C_c^\infty(\C)$,
\[
\Big\langle \frac{\lambda}{c_\delta}\psi_\delta\, e^{-\pi|\cdot|^2},\varphi\Big\rangle
=
\frac{\lambda}{c_\delta}\int_\C \psi_\delta(z)\,e^{-\pi|z|^2}\varphi(z)\,dA(z).
\]
Decompose $\varphi(z)e^{-\pi|z|^2}=\varphi(0)+r(z)$ with $r$ continuous and $r(0)=0$; then
\[
\frac{1}{c_\delta}\int_\C \psi_\delta(z)\,e^{-\pi|z|^2}\varphi(z)\,dA(z)
=
\varphi(0)\cdot\frac{1}{c_\delta}\int_\C\psi_\delta(z)e^{-\pi|z|^2}\,dA(z)
+
\frac{1}{c_\delta}\int_\C\psi_\delta(z)\,r(z)\,dA(z).
\]
The first summand equals $\varphi(0)$. The second is bounded in absolute value by $c_\delta^{-1}\,\|\psi_\delta\|_{L^1}\,\sup_{|z|\le 2\delta}|r(z)|=c_\delta^{-1}\sup_{|z|\le 2\delta}|r(z)|\to 0$ as $\delta\to 0$, since $r$ is continuous, $r(0)=0$, and $c_\delta\to 1$. Multiplying by $\lambda$ proves \eqref{eq:psi-delta-to-mass}.

In particular, $\Delta\widetilde u_0\to\Delta u_*$ in $\mathcal D'(\C)$. To upgrade this to local uniform convergence, we use the Newtonian-potential representations of $\widetilde u_0$ and $u_*$ together with the harmonic correction enforcing compact support: writing
\[
\widetilde u_0(z)
=
\frac{1}{2\pi}\int_\C \log\frac{1}{|z-\zeta|}\Big(-e^{-\pi|\zeta|^2}\mathbf 1_{D_0}(\zeta)+\frac{\lambda}{c_\delta}\psi_\delta(\zeta)e^{-\pi|\zeta|^2}\Big)dA(\zeta)+h_\delta(z),
\]
\[
u_*(z)
=
\frac{1}{2\pi}\int_\C \log\frac{1}{|z-\zeta|}\Big(-e^{-\pi|\zeta|^2}\mathbf 1_{D_0}(\zeta)\Big)dA(\zeta)+\frac{\lambda}{2\pi}\log\frac{1}{|z|}+h_*(z),
\]
where $h_\delta,h_*$ are harmonic functions chosen so as to enforce compact support, the contribution from $-e^{-\pi|\zeta|^2}\mathbf 1_{D_0}$ cancels in the difference and one is left with
\[
\widetilde u_0(z)-u_*(z)
=
\frac{\lambda}{2\pi}\Big(\int_\C \log\frac{1}{|z-\zeta|}\frac{1}{c_\delta}\psi_\delta(\zeta)e^{-\pi|\zeta|^2}\,dA(\zeta)-\log\frac{1}{|z|}\Big)+\big(h_\delta(z)-h_*(z)\big).
\]
By dominated convergence applied to the logarithmic potential with kernel uniformly bounded on compact subsets of $\overline D_0\setminus\{0\}$, together with \eqref{eq:psi-delta-to-mass}, the displayed difference tends to zero locally uniformly on $\overline D_0\setminus\{0\}$. Combined with interior elliptic regularity for $\Delta(\widetilde u_0-u_*)$, which is harmonic on any compact subset of $D_0\setminus\{0\}$ once $\delta$ is so small that $\supp\psi_\delta$ is contained in any prescribed neighborhood of the origin, this gives
\[
\widetilde u_0\xrightarrow[\delta\to 0]{}u_* \qquad \text{locally uniformly on $\overline D_0\setminus\{0\}$}.
\]
Since $u_*$ is strictly positive on every compact subset of $D_0\setminus\{0\}$ and continuous up to $\partial D_0$ where it vanishes, this implies $\widetilde u_0>0$ on every fixed compact subset of $D_0\setminus\{0\}$ for all $\delta$ small enough; the strict inequality $u_*>0$ at interior points provides the room needed to absorb the local-uniform error. On the small annulus $\delta\le|z|\le 2\delta$, where $\psi_\delta$ is supported, the explicit profile near the origin shows that $\widetilde u_0$ is comparable to a truncated logarithm $-\log(|z|/r_0)$ shifted by an $O(1)$ smooth correction, and is therefore strictly positive. We conclude that $\widetilde u_0>0$ throughout $D_0$ for $\delta>0$ small.

\smallskip
\emph{Perturbed problem via Isakov.}
We now invoke the perturbative free-boundary theorem of Isakov \cite{Isakov}, given by Proposition~\ref{prop:isakov-tailored}, verifying its hypotheses carefully. Set
\[
f_0(z):=e^{-\pi|z|^2}\Big(1-\frac{\lambda}{c_\delta}\psi_\delta(z)\Big).
\]
Since $\psi_\delta\in C_c^\infty(\C)$, we have $f_0\in C^\infty(\R^2)$, hence in particular $f_0\in C^5(\R^2)$. If $\delta$ is chosen so small that $\supp\psi_\delta\subset D_0$, then on $\partial D_0$,
\[
f_0(z)=e^{-\pi|z|^2}>0.
\]
Then \eqref{eq:unperturbed-free-boundary-new} becomes
\[
\Delta \widetilde u_0=f_0\,\mathbf{1}_{D_0}
\qquad \text{in }\C,
\]
with \(\widetilde u_0=0\) on \(\C\setminus D_0\) and \(\widetilde u_0>0\) in \(D_0\), by the positivity established above. Since \(\widetilde u_0\) was constructed as the compactly supported radial solution of \eqref{eq:unperturbed-free-boundary-new}, it is \(C^1\) across \(\partial D_0\); equivalently \(\widetilde u_0=\partial_\nu \widetilde u_0=0\) on \(\partial D_0\). Thus all hypotheses of Proposition~\ref{prop:isakov-tailored} are satisfied with \(\Omega_0=D_0\).

For the perturbed source, define
\[
f(z):=e^{-\pi|z|^2}\Big(|P(z)|^2-\frac{\lambda}{c_\delta}\psi_\delta(z)\Big).
\]
Since $P$ is a polynomial, $f\in C^\infty(\R^2)$ as well. Moreover, for any fixed disk $\Omega'\supset D_0$, the $C^5(\Omega')$-norm of $f-f_0$ is controlled by the coefficients of $P-1$: writing
\[
|P(z)|^2-1=\sum_{k=1}^N b_k z^k+\sum_{k=1}^N \overline{b_k}\,\overline z^{\,k}+\sum_{k,\ell=1}^N b_k\overline{b_\ell}\,z^k\overline z^{\,\ell},
\]
which is a polynomial in $z,\overline z$ that vanishes when all $b_k=0$ and depends polynomially on the coefficients, a standard estimate yields
\[
\|f-f_0\|_{C^5(\Omega')}\le C(\Omega',N,\lambda)\max_{1\le k\le N}|b_k|.
\]
Hence, if the coefficients $b_k$ are sufficiently small, we may arrange that
\[
\|f-f_0\|_{C^5(\Omega')}<\varepsilon_1,
\]
where $\varepsilon_1$ is the smallness threshold furnished by Proposition~\ref{prop:isakov-tailored}. Applying that proposition, we obtain a \(C^2\) domain \(U_0\), represented as a small normal deformation of \(D_0\), and a compactly supported function \(u_0\) such that
\[
\Delta u_0=f\,\mathbf{1}_{U_0}
\qquad \text{in }\C,
\]
with \(u_0=0\) on \(\C\setminus U_0\), \(\partial_\nu u_0=0\) on \(\partial U_0\), and \(u_0>0\) in \(U_0\).
If the perturbation is sufficiently small, then $U_0$ is a small $C^2$ deformation of $D_0$, and since $\supp\psi_\delta\subset D_0$ we still have
\[
\supp\psi_\delta\subset U_0.
\]
Therefore
\[
\Delta u_0
=
|P(z)|^2e^{-\pi|z|^2}\mathbf{1}_{U_0}
-\frac{\lambda}{c_\delta}\psi_\delta(z)e^{-\pi|z|^2},
\]
so \eqref{eq:inverse-first-new} holds. This proves \eqref{eq:first-domain-new}, and completes Step~1: we have produced a $C^2$ domain $U_0$, lying in a small $C^2$-neighborhood of the Daubechies disk $D_0$, on which the symmetrized moment identity holds with constant right-hand side $\lambda$. We will use $U_0$ in Step~2 as the unperturbed reference for the application of Proposition~\ref{prop:cherednichenko-tailored}.

\medskip
\noindent\textbf{Step 2: construction of a domain satisfying \eqref{eq:almost-reduced-new}.}
With the auxiliary domain $U_0$ at our disposal, we now seek a $C^2$ domain $U$ for which the right-hand side of the moment identity is the polynomial $\lambda Q(w)$ rather than the constant $\lambda$. The crucial observation is that integrating the parameter $w$ along the positive real axis against an exponential weight converts \eqref{eq:almost-reduced-new} into an explicit identity for a weighted Cauchy transform; this is exactly the form to which Proposition~\ref{prop:cherednichenko-tailored} of Cherednichenko \cite{Cherednichenko} applies. The reformulation we are about to perform is closely related to the quadrature-domain framework developed in \cite{Gustafsson,Gustafsson-Shahgholian,Karp,Shahgholian}.

Restrict \eqref{eq:almost-reduced-new} to real values $w=t>0$, multiply by $e^{-\pi \zeta t}$, and integrate in $t>0$. If
\[
\zeta>\sup_{z\in U}|z|,
\]
the integrand of the resulting double integral is absolutely integrable: indeed, for real $\zeta>\sup_{z\in U}|z|$,
\[
\Re(\overline z-\zeta)\le |z|-\zeta<0
\qquad \text{for every } z\in U,
\]
so
\[
\int_U \int_0^\infty \big|e^{\pi(\overline z-\zeta)t}\big|\,dt\,|P(z)|^2 e^{-\pi|z|^2}\,dA(z)
\le
\int_U \frac{|P(z)|^2 e^{-\pi|z|^2}}{\pi(\zeta-|z|)}\,dA(z)<\infty,
\]
because $U$ is bounded. Fubini's theorem therefore applies and yields
\[
\int_U \left(\int_0^\infty e^{\pi(\overline z-\zeta)t}\,dt\right)|P(z)|^2e^{-\pi|z|^2}\,dA(z)
=
\lambda\int_0^\infty e^{-\pi\zeta t}Q(t)\,dt.
\]
Writing
\[
Q(w)=\alpha_0+\alpha_1 w+\cdots+\alpha_n w^n,
\]
we obtain
\[
\frac{1}{\pi}\int_U \frac{|P(z)|^2e^{-\pi|z|^2}}{\zeta-\overline z}\,dA(z)
=
\lambda \sum_{k=0}^n \alpha_k \int_0^\infty e^{-\pi \zeta t} t^k\,dt.
\]
Since
\[
\int_0^\infty e^{-\pi \zeta t} t^k\,dt
=
\frac{\Gamma(k+1)}{(\pi \zeta)^{k+1}},
\]
it follows that
\[
\frac{1}{\pi}\int_U \frac{|P(z)|^2e^{-\pi|z|^2}}{\zeta-\overline z}\,dA(z)
=
\lambda \sum_{k=0}^n \alpha_k \frac{\Gamma(k+1)}{\pi^{k+1}\zeta^{k+1}}.
\]

To put this into the form of a weighted Cauchy transform, write
\[
E^\ast:=\{\overline z:z\in E\}
\]
for reflection across the real axis, and let
\[
V:=U^\ast,
\qquad
\mu(w):=|P(\overline w)|^2 e^{-\pi|w|^2}.
\]
The change of variables $w=\overline z$ yields
\[
\int_U \frac{|P(z)|^2e^{-\pi|z|^2}}{\zeta-\overline z}\,dA(z)
=
\int_V \frac{\mu(w)}{\zeta-w}\,dA(w)
=
g_V(\zeta).
\]
Hence
\begin{equation}\label{eq:analytic-construct-new}
g_V(\zeta)
=
\lambda \sum_{k=0}^n \alpha_k \frac{\Gamma(k+1)}{\pi^k \zeta^{k+1}},
\qquad \zeta\in \R_+, \ \zeta\gg 1.
\end{equation}
Since both sides are holomorphic for large $\zeta$, analytic continuation gives the desired identity on $\C\setminus\overline V$.

Conversely, suppose a domain $V$ satisfies \eqref{eq:analytic-construct-new}. Expanding $1/(\zeta-w)=\sum_{k\ge 0}w^k/\zeta^{k+1}$ for $|\zeta|$ large and matching coefficients yields the moments
\[
\int_V w^k\mu(w)\,dA(w)=\lambda\,\alpha_k\,\frac{\Gamma(k+1)}{\pi^k},
\qquad 0\le k\le n,
\]
together with the vanishing of the corresponding moments for $k>n$ (since $Q$ has degree $n$). These moment identities recover \eqref{eq:almost-reduced-new} for the reflected domain \(U=V^\ast\) via the Laplace-transform calculation read in reverse, in conjunction with the identity theorem for entire functions. Thus the reflection \(V=U^\ast\) is merely a bookkeeping device that places the kernel in the form $1/(\zeta-w)$ required by the weighted Cauchy transform formalism.

Let $g$ denote the rational function on the right-hand side of \eqref{eq:analytic-construct-new}. We now verify the hypotheses of Proposition~\ref{prop:cherednichenko-tailored}. Since
\[
Q(w)=P^\#(\partial_w/\pi)\big(P(w)\big),
\]
and $P(0)=1$, the constant term of $Q$ is
\[
Q(0)=1+\sum_{k=1}^N \frac{k!}{\pi^k}|b_k|^2.
\]
In particular,
\[
Q(0)>0,
\]
so the leading term in the expansion of $g$ at infinity is
\[
\frac{\lambda Q(0)}{z},
\]
and consequently
\[
g(z)\to 0,
\qquad
z\,g(z)\to \lambda Q(0)>0
\qquad \text{as } z\to\infty.
\]
Thus \eqref{eq:appendix-g-hypotheses} holds, with $c_0=\lambda Q(0)$.

We next verify the rational-data smallness condition in Proposition~\ref{prop:cherednichenko-tailored}. The weighted Cauchy transform formulation is written for the reflected domain, so set
\[
V_0:=U_0^\ast.
\]
Since \(U_0\) was obtained in Step~1 as a small perturbation of the Daubechies disk \(D_0\), the reflected domain \(V_0\) is also a small perturbation of \(D_0\). We choose concentric disks \(V_1\) and \(V_2\), centered at the center of \(D_0\), with radii, say, \((1-10^{-3})\) and \((1+10^{-3})\) times the radius of \(D_0\), after shrinking the perturbation so that
\[
\overline{V_1}\subset V_0,\qquad \overline{V_0}\subset V_2 .
\]
By Step~1,
\[
\int_{U_0}|P(z)|^2e^{-\pi|z|^2}e^{\pi \overline zw}\,dA(z)=\lambda,
\qquad w\in \C.
\]
Repeating the Laplace-transform argument above, with the constant $\lambda$ in place of the polynomial $\lambda Q(w)$, gives
\[
g_{V_0}(z)=\frac{\lambda}{z}
\qquad \text{for } z\in \C\setminus \overline{V_0},
\]
where \(g_{V_0}\) denotes the weighted Cauchy transform with weight \(\mu(w)=|P(\overline w)|^2e^{-\pi|w|^2}\). Consequently, on \(\partial V_2\),
\[
g(z)-g_{V_0}(z)
=
\lambda \sum_{k=1}^n \alpha_k \frac{\Gamma(k+1)}{\pi^k z^{k+1}}
+
\lambda\,(Q(0)-1)\,\frac1z .
\]
The right-hand side belongs to the finite-dimensional real rational space
\[
\mathcal R_n:=\operatorname{span}_{\R}
\{z^{-1},z^{-2},iz^{-2},\ldots,z^{-(n+1)},iz^{-(n+1)}\}.
\]
Every element of \(\mathcal R_n\) has real \(z^{-1}\)-coefficient. It is
holomorphic on \(\C\setminus\overline{V_1}\), has only a pole at the origin,
and vanishes at infinity. Since all norms on the finite-dimensional
space \(\mathcal R_n\) are equivalent, there is a constant
\(C(V_1,V_2,\lambda,n)\) such that
\[
\big\|(g-g_{V_0})|_{\partial V_2}\big\|_{C^2(\partial V_2)}
\le
C(V_1,V_2,\lambda,n)
\Big(|Q(0)-1|+\sum_{k=1}^n |\alpha_k|\Big).
\]
Each coefficient \(\alpha_k\) of \(Q-1\) depends polynomially on the
coefficients of \(P-1\) and their conjugates, with no constant term, and
therefore tends to zero with \(\max_{1\le j\le N}|b_j|\). Hence, by choosing
the coefficients sufficiently small, we may ensure that
\[
(g-g_{V_0})\in\mathcal R_n,
\qquad
\big\|(g-g_{V_0})|_{\partial V_2}\big\|_{C^2(\partial V_2)}<\varepsilon_0,
\]
where \(\varepsilon_0\) is the smallness threshold provided by Proposition~\ref{prop:cherednichenko-tailored}. Applying that proposition, we obtain a \(C^2\) domain \(\widetilde U\) such that
\[
g_{\widetilde U}(z)=g(z),
\qquad z\in \C\setminus \overline{\widetilde U}.
\]
Setting \(U:=\widetilde U^\ast\), the converse part of the Laplace-transform argument gives \eqref{eq:almost-reduced-new}.

Finally, since \(\widetilde U\subset V_2\), the domain \(U=\widetilde U^\ast\) is contained in \(V_2^\ast=V_2\), and \(V_2\) is fixed once \(\lambda\) is fixed. By shrinking the admissible size of the coefficients of \(P-1\) once more we may ensure that
\[
P^\#(z)\neq 0
\qquad \text{for all } z\in V_2.
\]
In particular,
\[
P^\#(z)\neq 0
\qquad \text{for all } z\in U.
\]
This non-vanishing of $P^\#$ on $U$ will play no direct role until Step~3, where it is needed to push the zeros of the characteristic polynomial $\widetilde P^\#(\zeta)=P^\#(\zeta/\pi)$ to large modulus.

\medskip
\noindent\textbf{Step 3: recovery of the original eigenvalue equation.}
It remains to upgrade the symmetrized identity \eqref{eq:almost-reduced-new} to the original equation \eqref{eq:eigenfunction-general-new}. Define
\[
G(w):=\int_U P(z)e^{-\pi|z|^2}e^{\pi \overline z w}\,dA(z),
\qquad w\in\C.
\]
Then $G$ is entire (the integrand is entire in $w$, and the dominated-convergence justification of differentiation under the integral sign carried out in Step~0 applies verbatim with $P$ in place of $|P|^2$). By construction, \eqref{eq:almost-reduced-new} holds, so applying $P^\#(\partial_w/\pi)$ to $G$ yields, by the same argument as in Step~0,
\[
P^\#(\partial_w/\pi)G(w)=\int_U|P(z)|^2 e^{-\pi|z|^2}e^{\pi\overline z w}\,dA(z)=\lambda Q(w)=P^\#(\partial_w/\pi)(\lambda P)(w),
\qquad w\in \C.
\]
Hence
\[
H:=G-\lambda P
\]
is an entire function satisfying the constant-coefficient linear differential equation
\[
P^\#(\partial_w/\pi)H=0.
\]

Introduce the characteristic polynomial of this operator,
\[
\widetilde P^\#(\zeta):=P^\#(\zeta/\pi),
\]
and let $\zeta_1,\dots,\zeta_\ell$ be its zeros, with multiplicities $m_1,\dots,m_\ell$. By the structure theorem for entire solutions of constant-coefficient linear differential equations of finite exponential type (see e.g.\ \cite[Ch.\ XII]{Hormander}), every entire solution of
\[
P^\#(\partial_w/\pi)H=0
\]
of at most exponential type is of the form
\begin{equation}\label{eq:exp-decomposition-step3}
H(w)=\sum_{j=1}^\ell R_j(w)\,e^{\zeta_j w},
\end{equation}
where each $R_j$ is a polynomial of degree strictly less than $m_j$.

We will reach a contradiction by showing that all $R_j$ must vanish. First, since $P^\#$ has no zeros on $V_2$, and the zeros of $\widetilde P^\#(\zeta)=P^\#(\zeta/\pi)$ are obtained by scaling those of $P^\#$ by the factor $\pi$, by choosing the coefficients of $P-1$ sufficiently small we may guarantee that all zeros of $\widetilde P^\#$ stay at distance at least $10\sup_{z\in U}|z|$ from the origin:
\[
|\zeta_j|>10\sup_{z\in U}|z|=:10M,
\qquad 1\le j\le \ell.
\]
(Indeed, the roots of $\widetilde P^\#$ depend continuously on its coefficients, and at $b=0$ the polynomial $\widetilde P^\#$ degenerates to the constant $1$, which has no roots; for small $b$ the roots therefore lie in any prescribed neighborhood of infinity.)

Next, we establish an upper bound on $H$ from its definition. Since $P$ is a fixed polynomial of degree $\le N$ and $U$ is bounded with $M=\sup_{z\in U}|z|$,
\[
|G(w)|
\le
\int_U |P(z)|e^{-\pi|z|^2}e^{\pi |z||w|}\,dA(z)
\le
\Big(\sup_{z\in U}|P(z)|e^{-\pi|z|^2}\Big)\,|U|\,e^{\pi M|w|}
\le
C(1+|w|)^N e^{\pi M|w|},
\qquad w\in\C,
\]
where the first inequality uses $|e^{\pi\overline z w}|\le e^{\pi|z||w|}$ and in the last inequality we have absorbed $\sup_{z\in U}|P(z)|\le C_0(1+M)^N$ and the area $|U|$ into the constant $C$. Since every polynomial is dominated by $e^{\varepsilon |w|}$ for arbitrary $\varepsilon>0$, it follows that for every $\varepsilon>0$ there exists $C_\varepsilon>0$ such that
\[
|G(w)|\le C_\varepsilon e^{(\pi+\varepsilon)M|w|},
\qquad w\in \C.
\]
As $\lambda P(w)$ grows only polynomially in $|w|$, the same estimate holds for $H$:
\[
|H(w)|\le C'_\varepsilon e^{(\pi+\varepsilon)M|w|}.
\]
In particular, $H$ has finite exponential type at most $\pi M$, which justifies the use of \eqref{eq:exp-decomposition-step3}.

Suppose, for contradiction, that some $R_j$ is not identically zero. By a small generic rotation of the coordinate $w$ if necessary, we may assume that the zeros $\zeta_j$ are arranged so that there is a unique $\zeta_1$ realizing
\[
|\zeta_1|=\max_j |\zeta_j|
\qquad\text{and}\qquad
R_1\not\equiv 0,
\]
and that, for every other zero $\zeta_j$ with $j\ne 1$,
\[
\Re(\zeta_j\overline{\zeta_1})<|\zeta_1|^2.
\]
(The strict inequality $\Re(\zeta_j\overline{\zeta_1})\le |\zeta_j||\zeta_1|<|\zeta_1|^2$ holds automatically when $|\zeta_j|<|\zeta_1|$; ties at maximal modulus are broken by the rotation.) Then, along the ray
\[
w=t\,\overline{\zeta_1},
\qquad t>0,
\]
the term $R_1(w)e^{\zeta_1 w}$ has modulus $|R_1(t\overline{\zeta_1})|e^{t|\zeta_1|^2}$, while every other term $R_j(w)e^{\zeta_j w}$ has modulus bounded by a polynomial in $t$ times $e^{t\,\Re(\zeta_j\overline{\zeta_1})}$, with $\Re(\zeta_j\overline{\zeta_1})<|\zeta_1|^2$. Hence the dominant exponential term yields
\[
|H(t\overline{\zeta_1})|\ge c\,e^{t|\zeta_1|^2}
\]
for all sufficiently large $t$, with $c>0$.

Evaluating the upper bound on $H$ at $w=t\overline{\zeta_1}$ gives
\[
|H(t\overline{\zeta_1})|
\le
C'_\varepsilon e^{(\pi+\varepsilon)Mt|\zeta_1|}.
\]
Comparing with the lower bound $|H(t\overline{\zeta_1})|\ge c\,e^{t|\zeta_1|^2}$ and letting $t\to\infty$, we obtain
\[
|\zeta_1|^2\le (\pi+\varepsilon)M|\zeta_1|,
\qquad\text{i.e.,}\qquad
|\zeta_1|\le (\pi+\varepsilon)M.
\]
Since $\varepsilon>0$ is arbitrary, this gives
\[
|\zeta_1|\le \pi M,
\]
contradicting our choice $|\zeta_1|>10M$ (recall that $\pi<10$). Hence every $R_j$ vanishes identically, so
\[
H\equiv 0.
\]
Equivalently,
\[
G=\lambda P,
\]
which is exactly \eqref{eq:eigenfunction-general-new}. This proves that $P$, and therefore $f_0$, is an eigenfunction of the localization operator associated with the domain $U$ with eigenvalue $\lambda$.
\end{proof}
\fi

% ==============================================================
% Section 4
% ==============================================================
